(1) 变化磁场满足 $\oint\mathbf E\cdot d\mathbf l=-d\Phi_B/dt$。理想无限长轴向均匀系统具有绕轴旋转与沿轴平移对称性，感生电场只能为 $E_\phi(r,t)\hat\phi$，故半径 $r$ 的圆环满足 $2\pi rE_\phi=-d\Phi_B(r,t)/dt$。
(2) 内环磁通 $\Phi_B=\pi r^2B(t)$，所以 $E_\phi=-r\dot B/2$（$r<a$）。外环磁通仅由螺线管截面贡献 $\Phi_B=\pi a^2B(t)$，故 $E_\phi=-a^2\dot B/(2r)$（$r>a$）。两式在 $r=a$ 连续。
(3) 闭合圆周电场功为 $W=q\oint\mathbf E\cdot d\mathbf l=-q\pi\min(r^2,a^2)\dot B$，在磁场变化时一般非零。相反，静电场为单值标量势的梯度，闭合回路积分恒为零，所以感应涡旋电场不能完全写成纯静电保守场。
(4) 因为 $\dot B=-B_0\omega\sin\omega t$，感生场为 $E_\phi=(B_0\omega/2)\sin\omega t\min(r,a^2/r)$。振幅在螺线管内部随 $r$ 线性增大、外部按 $1/r$ 下降，轴心趋零。忽略迟滞效应须有 $\omega a/c\ll1$。
